\documentclass{amsart}
% For dealing with references we use the comment environment
\usepackage{verbatim}
\newenvironment{reference}{\comment}{\endcomment}
%\newenvironment{reference}{}{}
\newenvironment{slogan}{\comment}{\endcomment}
\newenvironment{history}{\comment}{\endcomment}

% For commutative diagrams we use Xy-pic
\usepackage[all]{xy}

% We use 2cell for 2-commutative diagrams.
\xyoption{2cell}
\UseAllTwocells

% We use multicol for the list of chapters between chapters
\usepackage{multicol}

% This is generally recommended for better output
\usepackage{lmodern}
\usepackage[T1]{fontenc}

% For cross-file-references

% Package for hypertext links:
\usepackage{hyperref}

% For any local file, say "hello.tex" you want to link to please

% Theorem environments.
%
\theoremstyle{plain}
\newtheorem{theorem}[subsection]{Theorem}
\newtheorem{proposition}[subsection]{Proposition}
\newtheorem{lemma}[subsection]{Lemma}

\theoremstyle{definition}
\newtheorem{definition}[subsection]{Definition}
\newtheorem{example}[subsection]{Example}
\newtheorem{exercise}[subsection]{Exercise}
\newtheorem{situation}[subsection]{Situation}

\theoremstyle{remark}
\newtheorem{remark}[subsection]{Remark}
\newtheorem{remarks}[subsection]{Remarks}

\numberwithin{equation}{subsection}

% Macros
%
\def\lim{\mathop{\mathrm{lim}}\nolimits}
\def\colim{\mathop{\mathrm{colim}}\nolimits}
\def\Spec{\mathop{\mathrm{Spec}}}
\def\Hom{\mathop{\mathrm{Hom}}\nolimits}
\def\Ext{\mathop{\mathrm{Ext}}\nolimits}
\def\SheafHom{\mathop{\mathcal{H}\!\mathit{om}}\nolimits}
\def\SheafExt{\mathop{\mathcal{E}\!\mathit{xt}}\nolimits}
\def\Sch{\mathit{Sch}}
\def\Mor{\mathop{\mathrm{Mor}}\nolimits}
\def\Ob{\mathop{\mathrm{Ob}}\nolimits}
\def\Sh{\mathop{\mathit{Sh}}\nolimits}
\def\NL{\mathop{N\!L}\nolimits}
\def\CH{\mathop{\mathrm{CH}}\nolimits}
\def\proetale{{pro\text{-}\acute{e}tale}}
\def\etale{{\acute{e}tale}}
\def\QCoh{\mathit{QCoh}}
\def\Ker{\mathop{\mathrm{Ker}}}
\def\Im{\mathop{\mathrm{Im}}}
\def\Coker{\mathop{\mathrm{Coker}}}
\def\Coim{\mathop{\mathrm{Coim}}}

% Boxtimes
%
\DeclareMathSymbol{\boxtimes}{\mathbin}{AMSa}{"02}

%
% Macros for moduli stacks/spaces
%
\def\QCohstack{\mathcal{QC}\!\mathit{oh}}
\def\Cohstack{\mathcal{C}\!\mathit{oh}}
\def\Spacesstack{\mathcal{S}\!\mathit{paces}}
\def\Quotfunctor{\mathrm{Quot}}
\def\Hilbfunctor{\mathrm{Hilb}}
\def\Curvesstack{\mathcal{C}\!\mathit{urves}}
\def\Polarizedstack{\mathcal{P}\!\mathit{olarized}}
\def\Complexesstack{\mathcal{C}\!\mathit{omplexes}}
% \Pic is the operator that assigns to X its picard group, usage \Pic(X)
% \Picardstack_{X/B} denotes the Picard stack of X over B
% \Picardfunctor_{X/B} denotes the Picard functor of X over B
\def\Pic{\mathop{\mathrm{Pic}}\nolimits}
\def\Picardstack{\mathcal{P}\!\mathit{ic}}
\def\Picardfunctor{\mathrm{Pic}}
\def\Deformationcategory{\mathcal{D}\!\mathit{ef}}

\setlength{\emergencystretch}{3em}
\begin{document}
\title[Stacks Project, Tag 089A]{377 --- Grothendieck algebraization of proper formal schemes with a compatible ample line bundle}
\maketitle
\noindent Copyright (C) 2005--2025 Johan de Jong. Stacks Project authors. Source: \href{https://stacks.math.columbia.edu/tag/089A}{Tag 089A}.

\noindent Editorial difficulty note (not author text). Difficulty H3: The proof must simultaneously lift a projective embedding through all infinitesimal levels with a uniform choice of degree. It builds a finite graded algebra from successive ideal layers, uses uniform H1 vanishing to lift the embedding sections, and algebraizes the resulting compatible system of closed projective subschemes..

\noindent Editorial provenance note (not author text). History: excerpt from revision \texttt{a0444\allowbreak 6e57e\allowbreak c1fbc\allowbreak 252a8\allowbreak 71afc\allowbreak ec775\allowbreak 2fb28\allowbreak 07b14}, coherent.tex, lines 7979--8115. Reference presentation was adapted; original-author context and core support blocks were retained or added. The mathematical wording is unchanged. Licensed under GNU FDL 1.2 or later; no invariant sections or cover texts. See ../licenses/COPYING. Context blocks reproduce author definitions and supporting results; they are not additional counted problems.

\begin{remark}[Unwinding Grothendieck's existence theorem]
\label{remark-reformulate-existence-theorem}
Let $A$ be a Noetherian ring complete with respect to an ideal $I$.
Write $S = \Spec(A)$ and $S_n = \Spec(A/I^n)$.
Let $X \to S$ be a separated morphism of finite type.
For $n \geq 1$ we set $X_n = X \times_S S_n$.
Picture:
$$
\xymatrix{
X_1 \ar[r]_{i_1} \ar[d] & X_2 \ar[r]_{i_2} \ar[d] & X_3 \ar[r] \ar[d] &
\ldots & X \ar[d] \\
S_1 \ar[r] & S_2 \ar[r] & S_3 \ar[r] & \ldots & S
}
$$
In this situation we consider systems $(\mathcal{F}_n, \varphi_n)$
where
\begin{enumerate}
\item $\mathcal{F}_n$ is a coherent $\mathcal{O}_{X_n}$-module,
\item $\varphi_n : i_n^*\mathcal{F}_{n + 1} \to \mathcal{F}_n$
is an isomorphism, and
\item $\text{Supp}(\mathcal{F}_1)$ is proper over $S_1$.
\end{enumerate}
Theorem \href{https://stacks.math.columbia.edu/tag/088E}{theorem grothendieck existence (Tag 088E)} says that the
completion functor
$$
\begin{matrix}
\text{coherent }\mathcal{O}_X\text{-modules }\mathcal{F} \\
\text{with support proper over }A
\end{matrix}
\quad
\longrightarrow
\quad
\begin{matrix}
\text{systems }(\mathcal{F}_n) \\
\text{as above}
\end{matrix}
$$
is an equivalence of categories. In the special case that $X$ is
proper over $A$ we can omit the conditions on the supports.
\end{remark}

\begin{lemma}
\label{lemma-graded-finiteness}
Let $A$ be a Noetherian ring.
Let $B$ be a finitely generated graded $A$-algebra.
Let $f : X \to \Spec(A)$ be a proper morphism.
Set $\mathcal{B} = f^*\widetilde B$.
Let $\mathcal{F}$ be a quasi-coherent
graded $\mathcal{B}$-module of finite type.
\begin{enumerate}
\item For every $p \geq 0$ the graded $B$-module $H^p(X, \mathcal{F})$
is a finite $B$-module.
\item If $\mathcal{L}$ is an ample invertible $\mathcal{O}_X$-module,
then there exists an integer $d_0$ such that
$H^p(X, \mathcal{F} \otimes \mathcal{L}^{\otimes d}) = 0$
for all $p > 0$ and $d \geq d_0$.
\end{enumerate}
\end{lemma}

\begin{proof}
To prove this we consider the fibre product diagram
$$
\xymatrix{
X' = \Spec(B) \times_{\Spec(A)} X
\ar[r]_-\pi \ar[d]_{f'} &
X \ar[d]^f \\
\Spec(B) \ar[r] &
\Spec(A)
}
$$
Note that $f'$ is a proper morphism, see
Morphisms, Lemma \href{https://stacks.math.columbia.edu/tag/01W4}{lemma base change proper (Tag 01W4)}.
Also, $B$ is a finitely generated $A$-algebra, and hence
Noetherian (Algebra, Lemma \href{https://stacks.math.columbia.edu/tag/00FN}{lemma Noetherian permanence (Tag 00FN)}).
This implies that $X'$ is a Noetherian scheme
(Morphisms, Lemma \href{https://stacks.math.columbia.edu/tag/01T6}{lemma finite type noetherian (Tag 01T6)}).
Note that $X'$ is the relative spectrum of the quasi-coherent
$\mathcal{O}_X$-algebra $\mathcal{B}$ by
Constructions, Lemma \href{https://stacks.math.columbia.edu/tag/01LX}{lemma spec properties (Tag 01LX)}.
Since $\mathcal{F}$ is a quasi-coherent $\mathcal{B}$-module
we see that there is a unique quasi-coherent
$\mathcal{O}_{X'}$-module $\mathcal{F}'$ such that
$\pi_*\mathcal{F}' = \mathcal{F}$, see
Morphisms, Lemma \href{https://stacks.math.columbia.edu/tag/01SB}{lemma affine equivalence modules (Tag 01SB)}
Since $\mathcal{F}$ is finite type as a $\mathcal{B}$-module we
conclude that $\mathcal{F}'$ is a finite type
$\mathcal{O}_{X'}$-module (details omitted). In other words,
$\mathcal{F}'$ is a coherent $\mathcal{O}_{X'}$-module
(Lemma \href{https://stacks.math.columbia.edu/tag/01XZ}{lemma coherent Noetherian (Tag 01XZ)}).
Since the morphism $\pi : X' \to X$ is affine we have
$$
H^p(X, \mathcal{F}) = H^p(X', \mathcal{F}')
$$
by Lemma \href{https://stacks.math.columbia.edu/tag/089W}{lemma relative affine cohomology (Tag 089W)}.
Thus (1) follows from
Lemma \href{https://stacks.math.columbia.edu/tag/02O6}{lemma proper over affine cohomology finite (Tag 02O6)}.
Given $\mathcal{L}$ as in (2) we set
$\mathcal{L}' = \pi^*\mathcal{L}$. Note that $\mathcal{L}'$ is
ample on $X'$ by
Morphisms, Lemma \href{https://stacks.math.columbia.edu/tag/0892}{lemma pullback ample tensor relatively ample (Tag 0892)}.
By the projection formula
(Cohomology, Lemma \href{https://stacks.math.columbia.edu/tag/01E8}{lemma projection formula (Tag 01E8)}) we have
$\pi_*(\mathcal{F}' \otimes \mathcal{L}') = \mathcal{F} \otimes \mathcal{L}$.
Thus part (2) follows by the same reasoning as above from
Lemma \href{https://stacks.math.columbia.edu/tag/02O1}{lemma kill by twisting (Tag 02O1)}.
\end{proof}


\begin{lemma}
\label{lemma-algebraize-formal-closed-subscheme}
Let $A$ be a Noetherian ring complete with respect to an ideal $I$.
Write $S = \Spec(A)$ and $S_n = \Spec(A/I^n)$.
Let $X \to S$ be a separated morphism of finite type.
For $n \geq 1$ we set $X_n = X \times_S S_n$.
Suppose given a commutative diagram
$$
\xymatrix{
Z_1 \ar[r] \ar[d] & Z_2 \ar[r] \ar[d] & Z_3 \ar[r] \ar[d] & \ldots \\
X_1 \ar[r]^{i_1} & X_2 \ar[r]^{i_2} & X_3 \ar[r] & \ldots
}
$$
of schemes with cartesian squares. Assume that
\begin{enumerate}
\item $Z_1 \to X_1$ is a closed immersion, and
\item $Z_1 \to S_1$ is proper.
\end{enumerate}
Then there exists a closed immersion of schemes $Z \to X$ such that
$Z_n = Z \times_S S_n$. Moreover, $Z$ is proper over $S$.
\end{lemma}

\begin{proof}
Let's write $j_n : Z_n \to X_n$ for the vertical morphisms.
As the squares in the statement are cartesian
we see that the base change of $j_n$ to $X_1$ is $j_1$.
Thus Morphisms, Lemma \href{https://stacks.math.columbia.edu/tag/0896}{lemma check closed infinitesimally (Tag 0896)}
shows that $j_n$ is a closed immersion.
Set $\mathcal{F}_n = j_{n, *}\mathcal{O}_{Z_n}$, so that
$j_n^\sharp$ is a surjection $\mathcal{O}_{X_n} \to \mathcal{F}_n$.
Again using that the squares are cartesian we see that
the pullback of $\mathcal{F}_{n + 1}$ to $X_n$ is $\mathcal{F}_n$.
Hence Grothendieck's existence theorem, as reformulated in
Remark \ref{remark-reformulate-existence-theorem},
tells us there exists a map
$\mathcal{O}_X \to \mathcal{F}$
of coherent $\mathcal{O}_X$-modules whose restriction to
$X_n$ recovers $\mathcal{O}_{X_n} \to \mathcal{F}_n$.
Moreover, the support of $\mathcal{F}$ is proper over $S$.
As the completion functor is exact (Lemma \href{https://stacks.math.columbia.edu/tag/0881}{lemma exact (Tag 0881)})
we see that the cokernel $\mathcal{Q}$ of $\mathcal{O}_X \to \mathcal{F}$
has vanishing completion. Since $\mathcal{F}$ has support
proper over $S$ and so does $\mathcal{Q}$ this implies that
$\mathcal{Q} = 0$ for example because the functor
(\href{https://stacks.math.columbia.edu/tag/088D}{equation completion functor proper over A (Tag 088D)}) is an equivalence
by Grothendieck's existence theorem.
Thus $\mathcal{F} = \mathcal{O}_X/\mathcal{J}$
for some quasi-coherent sheaf of ideals $\mathcal{J}$.
Setting $Z = V(\mathcal{J})$ finishes the proof.
\end{proof}


\begin{theorem}[Grothendieck's algebraization theorem]
\label{theorem-algebraization}
Let $A$ be a Noetherian ring complete with respect to an ideal $I$.
Set $S = \Spec(A)$ and $S_n = \Spec(A/I^n)$. Consider a commutative
diagram
$$
\xymatrix{
X_1 \ar[r]_{i_1} \ar[d] & X_2 \ar[r]_{i_2} \ar[d] & X_3 \ar[r] \ar[d] &
\ldots \\
S_1 \ar[r] & S_2 \ar[r] & S_3 \ar[r] & \ldots
}
$$
of schemes with cartesian squares. Suppose given $(\mathcal{L}_n, \varphi_n)$
where each $\mathcal{L}_n$ is an invertible sheaf on $X_n$ and
$\varphi_n : i_n^*\mathcal{L}_{n + 1} \to \mathcal{L}_n$
is an isomorphism. If
\begin{enumerate}
\item $X_1 \to S_1$ is proper, and
\item $\mathcal{L}_1$ is ample on $X_1$
\end{enumerate}
then there exists a proper morphism of schemes $X \to S$
and an ample invertible $\mathcal{O}_X$-module $\mathcal{L}$
and isomorphisms $X_n \cong X \times_S S_n$ and
$\mathcal{L}_n \cong \mathcal{L}|_{X_n}$ compatible with
the morphisms $i_n$ and $\varphi_n$.
\end{theorem}

\begin{proof}
Since the squares in the diagram are cartesian and since the morphisms
$S_n \to S_{n + 1}$ are closed immersions, we see that the morphisms
$i_n$ are closed immersions too. In particular we may think of
$X_m$ as a closed subscheme of $X_n$ for $m < n$. In fact $X_m$ is
the closed subscheme cut out by the quasi-coherent sheaf of ideals
$I^m\mathcal{O}_{X_n}$. Moreover, the underlying topological spaces
of the schemes $X_1, X_2, X_3, \ldots$ are all identified, hence we
may (and do) think of sheaves $\mathcal{O}_{X_n}$ as living on the
same underlying topological space; similarly for coherent
$\mathcal{O}_{X_n}$-modules. Set
$$
\mathcal{F}_n =
\Ker(\mathcal{O}_{X_{n + 1}} \to \mathcal{O}_{X_n})
$$
so that we obtain short exact sequences
$$
0 \to \mathcal{F}_n \to \mathcal{O}_{X_{n + 1}} \to \mathcal{O}_{X_n} \to 0
$$
By the above we have $\mathcal{F}_n = I^n\mathcal{O}_{X_{n + 1}}$.
It follows $\mathcal{F}_n$ is a coherent sheaf on $X_{n + 1}$
annihilated by $I$, hence we may (and do) think of it as a coherent
module $\mathcal{O}_{X_1}$-module. Observe that for $m > n$ the sheaf
$$
I^n\mathcal{O}_{X_m}/I^{n + 1}\mathcal{O}_{X_m}
$$
maps isomorphically to $\mathcal{F}_n$ under the map
$\mathcal{O}_{X_m} \to \mathcal{O}_{X_{n + 1}}$. Hence given
$n_1, n_2 \geq 0$ we can pick an $m > n_1 + n_2$ and consider the
multiplication map
$$
I^{n_1}\mathcal{O}_{X_m} \times I^{n_2}\mathcal{O}_{X_m}
\longrightarrow
I^{n_1 + n_2}\mathcal{O}_{X_m} \to \mathcal{F}_{n_1 + n_2}
$$
This induces an $\mathcal{O}_{X_1}$-bilinear map
$$
\mathcal{F}_{n_1} \times \mathcal{F}_{n_2} \longrightarrow
\mathcal{F}_{n_1 + n_2}
$$
which in turn defines the structure of a graded $\mathcal{O}_{X_1}$-algebra
on $\mathcal{F} = \bigoplus_{n \geq 0} \mathcal{F}_n$.

\medskip\noindent
Set $B = \bigoplus I^n/I^{n + 1}$; this is a finitely generated
graded $A/I$-algebra. Set $\mathcal{B} = (X_1 \to S_1)^*\widetilde{B}$.
The discussion above provides us with a canonical surjection
$$
\mathcal{B} \longrightarrow \mathcal{F}
$$
of graded $\mathcal{O}_{X_1}$-algebras. In particular we see that
$\mathcal{F}$ is a finite type quasi-coherent graded $\mathcal{B}$-module.
By Lemma \ref{lemma-graded-finiteness} we can find an integer $d_0$
such that $H^1(X_1, \mathcal{F} \otimes \mathcal{L}^{\otimes d}) = 0$
for all $d \geq d_0$. Pick a $d \geq d_0$ such that there exist sections
$s_{0, 1}, \ldots, s_{N, 1} \in \Gamma(X_1, \mathcal{L}_1^{\otimes d})$
which induce an immersion
$$
\psi_1 : X_1 \to \mathbf{P}^N_{S_1}
$$
over $S_1$, see
Morphisms, Lemma \href{https://stacks.math.columbia.edu/tag/01VT}{lemma finite type over affine ample very ample (Tag 01VT)}.
As $X_1$ is proper over $S_1$ we see that $\psi_1$
is a closed immersion, see
Morphisms, Lemma \href{https://stacks.math.columbia.edu/tag/01W6}{lemma image proper scheme closed (Tag 01W6)}
and
Schemes, Lemma \href{https://stacks.math.columbia.edu/tag/01IQ}{lemma immersion when closed (Tag 01IQ)}.
We are going to ``lift'' $\psi_1$ to a compatible system of
closed immersions of $X_n$ into $\mathbf{P}^N$.

\medskip\noindent
Upon tensoring the short exact sequences of the first paragraph
of the proof
by $\mathcal{L}_{n + 1}^{\otimes d}$ we obtain short exact sequences
$$
0 \to \mathcal{F}_n \otimes \mathcal{L}_{n + 1}^{\otimes d} \to
\mathcal{L}_{n + 1}^{\otimes d} \to \mathcal{L}_{n + 1}^{\otimes d} \to 0
$$
Using the isomorphisms $\varphi_n$ we obtain isomorphisms
$\mathcal{L}_{n + 1} \otimes \mathcal{O}_{X_l} = \mathcal{L}_l$
for $l \leq n$. Whence the sequence above becomes
$$
0 \to \mathcal{F}_n \otimes \mathcal{L}_1^{\otimes d} \to
\mathcal{L}_{n + 1}^{\otimes d} \to \mathcal{L}_n^{\otimes d} \to 0
$$
The vanishing of $H^1(X, \mathcal{F}_n \otimes \mathcal{L}_1^{\otimes d})$
implies we can inductively lift
$s_{0, 1}, \ldots, s_{N, 1} \in \Gamma(X_1, \mathcal{L}_1^{\otimes d})$
to sections
$s_{0, n}, \ldots, s_{N, n} \in \Gamma(X_n, \mathcal{L}_n^{\otimes d})$.
Thus we obtain a commutative diagram
$$
\xymatrix{
X_1 \ar[r]_{i_1} \ar[d]_{\psi_1} &
X_2 \ar[r]_{i_2} \ar[d]_{\psi_2} &
X_3 \ar[r] \ar[d]_{\psi_3} &
\ldots \\
\mathbf{P}^N_{S_1} \ar[r] &
\mathbf{P}^N_{S_2} \ar[r] &
\mathbf{P}^N_{S_3} \ar[r] & \ldots
}
$$
where
$\psi_n = \varphi_{(\mathcal{L}_n, (s_{0, n}, \ldots, s_{N, n}))}$
in the notation of Constructions, Section
\href{https://stacks.math.columbia.edu/tag/01ND}{section projective space (Tag 01ND)}.
As the squares in the statement of the theorem are cartesian
we see that the squares in the above diagram are cartesian.
We win by applying Lemma \ref{lemma-algebraize-formal-closed-subscheme}.
\end{proof}
\end{document}
